% UTUgradu ja saavutettavuusohjeet täyttävä latex-pohja graduille
% 02/2021 Petriina Paturi

% UTUGradu and accessibility directives fulfilling latex-file for master's thesis
% Translated 02/2021 Christian Palmroos

% Metadata for the pdf/a file format. Remember to recompile!
\begin{filecontents*}{\jobname.xmpdata}
\Title{masters_thesis}
\Author{S. Student}
\end{filecontents*}

% Download utugradu2020.cls from intranet.
% pdfpagelabels makes the pdf have same page numbers that appear on the page.
% Accessibility requires this.
\documentclass[a4paper,pdfa,pdfpagelabels,colorlinks]{utugradu2020_en}  

\usepackage[utf8]{inputenc} %makes utf-8 work
\usepackage[T1]{fontenc}
\usepackage{tytiivis_en} %For the synopsis page; download tytiivis.sty from intranet
\usepackage{graphicx} %Enables figures
\usepackage{lastpage}
\usepackage[sort&compress, comma, square, numbers]{natbib} %for references
\usepackage{amssymb, amsthm, amsmath}
\usepackage{mathtools}

%\DeclareMathSymbol{.}{\mathpunct}{letters}{"3B}  % Dot for decimals
%\DeclareMathSymbol{.}{\mathord}{letters}{"3B}
%\DeclareMathSymbol{\decimal}{\mathord}{letters}{"3A}


%\usepackage{axessibility} % Accessibility package makes equations more readable

% Thesis pdf must be up to pdf/a-3b standard. This and pdflatex enable it.
\usepackage[a-3b]{pdfx}

% This is required for for archiving purposes.
\usepackage{hyperref}
% pdflatex requires figures to not be of eps-type. For example, pdf works well for vector graphics, as long as they are the same pdf/a-3b type, and png for pixel graphics.

% Make hyperlinks not visible in the printed version.
%\hypersetup{
%    colorlinks = false,
%    linkbordercolor = {white},
%    citebordercolor = {white}
%}
% With this there should be no white border around the links (the above works in the Finnish version 
% for some reason.
\hypersetup{
    %colorlinks = true,
    linkcolor= {black},
    citecolor = {black}
}

% Examples of defining new commands:
% Command \mathbi{``vector symbol''} creates bold italic letter. 
% For Greek alphabet you probably need to use \pmb
\newcommand{\mathbi}[1]{\textbf{\em #1}}



\begin{document}
\pagenumbering{roman} % roman numerals for pages until the primary text starts

\title{Title}                   % Title of thesis
\thesis{Master's thesis}        % Kind of thesis (Bachelor's, Master's)
\author{B.Sc. Susan Student}    % Author
\year{2021}                     % Year
\subject{Physics}               % Subject
\examiner{Prof. P.P.}           % Examiner #1
\examiner{Dr. H.H.}             % Examiner #2 (comment out if only one examiner)
%\examiner{Dr. A.A.}             % Examiner #3 (comment out if only 1-2 examiners)


\maketitle
\newpage
\thispagestyle{empty}
\vspace*{10cm}

\vfill

% Note that you will have to print out this page as two-sided with the cover of the thesis
\hspace*{-2cm}\parbox{\textwidth}{The originality of this thesis has been checked in accordance with the University of Turku quality assurance system using Turnitin Originality Check service.} 


\newpage
%This is the synopsis:
\begin{synopsis}
        {Department of Physics and Astronomy}
        {Student, Susan}
        {Title of thesis}
        {Master's thesis, \pageref{LastPage} pp.}
        {Physics}
        {April 2021}
	
	Synopsis here!
	
	Keywords: keyword1, keyword2
\end{synopsis}


\tableofcontents

\newpage
\pagenumbering{arabic} % Start numerating pages in Arabic numerals
\section*{Preface}
\addcontentsline{toc}{section}{Preface}
%Preface is not to be numerated in the table of contents


When writing your thesis, remember these basic rules:
\begin{enumerate}
\item All figures must be referenced in the text, e.g. "One can see from figure \ref{fig1}, that referencing to figures in latex is easy".
\item Figures and tables belong to the top of the page. Latex does this automatically.

\item Figures should be reasonably compact, so that its area is fully used.

\item Captions should include enough information to understand the respective figure/table without reading the text. Abbreviations and physical quantities are to be expressed. This improves accessibility. Latex does not support alternative texts as of now.

\item  Remember to reference when taking a figure from a source. Theses go to a public digital archive: remember copyright!

\item Introduce all abbreviations when first using them: e.g. astronomical unit (AU)

\item If you have to come up with an original translation for a word, add the original word after first usage of your translation. For example: "braking radiation" (Bremsstrahlung).

\item Physical quantities are written in italic, e.g $\rho = m/V$. Units are written in roman, e.g. 1 m$^2$. Vectors are written in bold italic, e.g. $\mathbi{v}$.

\item Equations are part of the text, commas and periods belong in the equations. Continue with a lower case letter after comma. Don't add an empty line after an equation.

\item Introduce all physical quantities after an equation. For example, Newton's second law is
\begin{equation}
\mathbi{F} = m\mathbi{a},
\end{equation}
where $\mathbi{F}$ is the force acting on an object, $m$ is the mass of an object, and $\mathbi{a}$ is the acceleration of that object. For Greek letter vectors, use the boldsymbol command: $\boldsymbol{\omega}$.

\item If all of the information in a paragraph comes from a singular source, then the citation belongs in the end of the paragraph, after the period. In all other cases the reference belongs before the period. Remember to cite your source every time you introduce numerical values or exact information to the text.

\end{enumerate}


% An example of how to produce a figure:
\begin{figure}
\begin{center}
\setlength{\unitlength}{1cm}
\begin{picture}(6,6)(-3,-3)
\put(-1.5,0){\vector(1,0){3}}
\put(2.7,-0.1){$\chi$}
\put(0,-1.5){\vector(0,1){3}}
\multiput(-2.5,1)(0.4,0){13}
{\line(1,0){0.2}}
\multiput(-2.5,-1)(0.4,0){13}
{\line(1,0){0.2}}
\put(0.2,1.4)
{$\beta=v/c=\tanh\chi$}
\qbezier(0,0)(0.8853,0.8853)
(2,0.9640)
\qbezier(0,0)(-0.8853,-0.8853)
(-2,-0.9640)
%\put(-3,-2){\circle*{0.2}}
\end{picture}

% This is how to import an image. If the picture is a pdf, it must be of pdf/a-3b type.
% includegraphics[width=10cm]{image.png}
\caption{Hyperbolic tangent function approaches asymptotically the values $\pm 1$}

\label{fig1}
\end{center}
\end{figure}

\section{Theory}

\newpage

% This makes line spacing smaller in references list. The font has to be changed for the command to work!
\renewcommand{\baselinestretch}{1}\large\normalsize

% Bibliography absolutely must be used.
\begin{thebibliography}{50}

\bibitem{lshort} T. Oetiker, H. Partl, I. Hyna and E. Schlegl,
Not so short introduction to \LaTeX 2e, 1998
\end{thebibliography}

% An alternative way for thebibliography environment is an BibTeX database, which you can create yourself,
% or use an existing (e.g. Wihuri's) database. We recommend this!

% BibTeX database can easily be created with TeXMaker. Then just recompile the .tex file, .bib file, and 
% once again .tex file. This should put your references in order.

%\bibliography{/var/bib/yhdistetty}
%\bibliographystyle{wihuri}

\end{document}
